\SourcePage{308}{311}
\section{Electron states in an atom}

An atom containing more than one electron is a complicated system of mutually
interacting electrons moving in the nuclear field. Strictly speaking, for such
a system one can consider only states of the system as a whole. Nevertheless,
it proves possible, to good accuracy, to introduce a state for each electron
separately: a stationary state of its motion in an effective spherically
symmetric field produced jointly by the nucleus and all the other electrons.
In general, these fields differ from one electron to another within the atom.
They must also all be determined simultaneously, since each depends on the
states of every other electron. A field of this kind is called
\emph{self-consistent}.

Since the self-consistent field has central symmetry, every electron state
has a definite value of its orbital angular momentum $l$.  For a fixed $l$,
the states of an individual electron are numbered in ascending order of energy
by the \emph{principal quantum number} $n$, whose possible values are
$n=l+1,l+2,\ldots$; this ordering convention is chosen to match the one used
for the hydrogen atom.  In a complex atom, however, the order in which levels
of different $l$ rise in energy generally differs from that in hydrogen.  For
hydrogen the energy is independent of $l$ altogether, and consequently states
with a larger $n$ always have a higher energy.  In a complex atom, by contrast,
the level with, for example, $n=5$, $l=0$ lies below the level with $n=4$,
$l=2$ (see \S~73 for further details).


An individual-electron state with specified $n$ and $l$ is conventionally
denoted by a symbol comprising a numeral for the principal quantum number and
a letter for the value of $l$\footnote{Another terminology refers to electrons
with principal quantum numbers $n=1,2,3,\ldots$ as electrons of the $K$,
$L$, $M$, \ldots{} shells, respectively (see \S~74).}. Thus $4d$ denotes the
state $n=4$, $l=2$. A complete description of an atomic state requires, in
addition to the values of the total $L,S,J$, a listing of the states occupied
by all its electrons. For example, the symbol $1s\,2p\,{}^3P_0$ designates a
helium-atom state with $L=1$, $S=1$, $J=0$, whose two electrons occupy

\SourcePage{309}{312}
the $1s$ and $2p$ states. When several electrons have the same $l$ and $n$,
this is abbreviated by an exponent; thus $3p^2$ denotes two electrons in
$3p$ states. The distribution of an atom's electrons among states with
different $l$ and $n$ is called its \emph{electron configuration}.

At fixed $n$ and $l$, an electron may have different projections of its
orbital angular momentum ($m$) and spin ($\sigma$) on the $z$ axis. For a
given $l$, the number $m$ takes $2l+1$ values, whereas $\sigma$ is restricted
to the two values $\pm1/2$. There are consequently $2(2l+1)$ different states
with the same $n,l$; these states are said to be \emph{equivalent}. By the
Pauli principle, each can contain one electron. Hence no more than $2(2l+1)$
electrons in an atom can simultaneously have identical $n,l$. A set of
electrons occupying every state with specified $n,l$ is called a
\emph{closed shell} of that type.

The difference between the energies of atomic levels having different $L,S$
but the same electron configuration\footnote{Here we disregard the fine
structure of each multiplet level.} results from the electrostatic interaction
of the electrons. Usually these energy differences are comparatively small,
being several times smaller than the separations between levels belonging to
different configurations. The relative ordering of levels with the same
configuration but different $L,S$ is described by the following empirically
established \emph{Hund rule} (F.~Hund, 1925):

\emph{The term of lowest energy is the one with the largest value of $S$
permitted by the electron configuration and, for that $S$, the largest
permitted value of $L$}\footnote{The requirement that $S$ be maximal can be
justified as follows. Consider, for example, a two-electron system. It may
have $S=0$ or $S=1$, and spin 1 corresponds to an antisymmetric coordinate
wave function $\varphi(\mathbf r_1,\mathbf r_2)$. Such a function vanishes
when $\mathbf r_1=\mathbf r_2$; in other words, in the $S=1$ state the
probability that both electrons are close together is small. Their
electrostatic repulsion is therefore comparatively weaker, and the energy is
lower. Likewise, in a system of several electrons the greatest spin
corresponds to the ``most antisymmetric'' coordinate wave function.}.

We now show how the atomic terms permitted for a specified electron
configuration can be found. If the electrons are not equivalent, the possible
values of $L,S$ follow directly from the rule for angular-momentum addition.
For instance, in the configuration $np,n'p$ (with different $n,n'$), the
total

\SourcePage{310}{313}
angular momentum $L$ can take the values 2, 1, 0, while the total spin is
$S=0,1$. Combining these values gives the terms ${}^{1/3}S$, ${}^{1/3}P$,
${}^{1/3}D$.

For equivalent electrons, however, the Pauli principle imposes restrictions.
As an example, consider a configuration of three equivalent $p$ electrons.
For $l=1$ (a $p$ state), the orbital-angular-momentum projection $m$ can be
$1,0,-1$, so there are six states with the following pairs $m,\sigma$:
\[
\begin{array}{lll}
\mathrm a)&1,,1/2,&\mathrm b)\ 0,,1/2\qquad \mathrm c)\ -1,,1/2,\\
\mathrm a')&1,,-1/2,&\mathrm b')\ 0,,-1/2\qquad \mathrm c')\ -1,,-1/2.
\end{array}
\]
The three electrons can be placed, one apiece, in any three of these states.
The resulting atomic states have the following values of the projections
$M_L=\sum m$, $M_S=\sum\sigma$ of the total orbital angular momentum and
spin:
\[
\begin{array}{lll}
\mathrm a+\mathrm a'+\mathrm b)&2,,1/2,&
\mathrm a+\mathrm a'+\mathrm c)\ 1,,1/2,
\quad \mathrm a+\mathrm b+\mathrm c)\ 0,,3/2,\\
&&\mathrm a+\mathrm b+\mathrm b')\ 1,,1/2,
\quad \mathrm a+\mathrm b+\mathrm c')\ 0,,1/2,\\
&&\mathrm a+\mathrm b'+\mathrm c)\ 0,,1/2,\\
&&\mathrm a'+\mathrm b+\mathrm c)\ 0,,1/2.
\end{array}
\]
(States with negative $M_L,M_S$ need not be written out, since they add
nothing new.) The presence of a state with $M_L=2$, $M_S=1/2$ shows that a
${}^2D$ term must occur; one state each with $(1,1/2)$ and $(0,1/2)$ must also
belong to it. One state with $(1,1/2)$ then remains, so there must be a
${}^2P$ term; one of the $(0,1/2)$ states belongs to it as well. Finally, the
remaining states $(0,3/2)$ and $(0,1/2)$ correspond to a ${}^4S$ term. Thus a
configuration of three equivalent $p$ electrons permits exactly one term of
each of the types ${}^2D$, ${}^2P$, and ${}^4S$.

Table~1 lists the possible terms for various configurations of equivalent
$p$- and $d$-electrons. A number written beneath a term symbol gives the
number of terms of that type in the configuration, whenever there is more
than one. A configuration containing the largest possible number of
equivalent electrons ($s^2,p^6,d^{10},\ldots$) always has the term ${}^1S$.
Notice that the sets of terms are the same for two configurations when one
contains as many electrons as the other lacks for a filled shell. This follows
directly because a missing electron in a shell may be regarded as a
\emph{hole}, whose state is specified by the same quantum numbers as the state
of the absent electron.

\SourcePage{311}{314}
When Hund's rule is used to determine the normal term of an atom from its
known electron configuration, only the unfilled shell need be considered,
since the angular momenta of the electrons in filled shells cancel one
another. Suppose, for example, that an atom has four
 $d$ electrons outside its closed
shells. The magnetic quantum number of a $d$ electron can take five values:
$0,\pm1,\pm2$. All four electrons can therefore have the same spin projection
$\sigma=1/2$, and the largest possible total spin is $S=2$. We must then
assign different values of $m$ to the electrons so as to maximize
$M_L=\sum m$; these are $2,1,0,-1$, giving $M_L=2$. It follows that the
greatest value of $L$ compatible with $S=2$ is likewise 2 (the ${}^5D$ term).

\begin{table}[ht]
\centering
\caption{All possible terms for configurations of equivalent electrons}
\begin{tabular}{c@{,\quad}c|lll}
$p$&$p^5$&${}^2P$&&\\
$p^2$&$p^4$&${}^1SD$&${}^3P$&\\
$p^3$&&${}^2PD$&${}^4S$&\\ \hline
$d$&$d^9$&${}^2D$&&\\
$d^2$&$d^8$&${}^1SDG$&${}^3PF$&\\
$d^3$&$d^7$&${}^2P\underset{2}{D}FGH$&${}^4PF$&\\
$d^4$&$d^6$&$\underset{2}{{}^1S}\underset{2}{D}F\underset{2}{G}I$
&$\underset{2}{{}^3P}D\underset{2}{F}GH$&${}^5D$\\
$d^5$&&${}^2SP\underset{3}{D}\underset{2}{F}\underset{2}{G}HI$
&${}^4P\underset{2}{D}\underset{2}{F}G$&${}^6S$
\end{tabular}
\end{table}

\ProblemHeading
Find the orbital wave functions for the possible states of a system of three
equivalent $p$-electrons.


\SolutionHeading In the ${}^4S$ state, the spin projections $\sigma$ of all
the electrons are equal, and hence their values of $m$ are distinct. The wave
function is given by a determinant of the form (61.5), constructed from the
functions $\psi_0$, $\psi_1$, and $\psi_{-1}$ (the subscript denotes $m$).

For the ${}^2D$ term, consider the state with the largest possible value
$M_L=2$. Two of the $m$ projections must then equal 1 and the remaining one
must equal 0. Let electrons 2 and 3 have $\sigma=+1/2$, and electron 1 have
$\sigma=-1/2$ (in accordance with the total spin $S=1/2$). The orbital wave
function with the required symmetry is
\[
\psi=\frac1{\sqrt2}\psi_1(1)
[\psi_0(2)\psi_1(3)-\psi_0(3)\psi_1(2)]
\]
(the numeral in a function's argument identifies the electron to which it
refers).

For the ${}^2P$ term, take the state with $M_L=1$ and the same electron-spin
projections as above. This state can be formed with

\SourcePage{312}{315}
two different sets of $m$ values, so its orbital wave function is the linear
combination
\[
\begin{gathered}
\psi=a\psi_{-111}+b\psi_{100},\\
\psi_{-111}=\psi_1(1)[\psi_{-1}(2)\psi_1(3)-\psi_{-1}(3)\psi_1(2)],\\
\psi_{100}=\psi_0(1)[\psi_1(2)\psi_0(3)-\psi_1(3)\psi_0(2)].
\end{gathered}
\]
To determine the coefficients, use the relation
\[
\hat L_+\psi=(\hat l_+^{(1)}+\hat l_+^{(2)}+\hat l_+^{(3)})\psi=0,
\]
which the wave function with $M_L=L$ must satisfy (see (27.8)). The matrix
elements (27.12) give
\[
\hat l_+\psi_1=0,\qquad \hat l_+\psi_{-1}=\sqrt2\psi_0,\qquad
\hat l_+\psi_0=\sqrt2\psi_1
\]
and hence
\[
\hat L_+\psi=\sqrt2(a-b)\psi_{011}=0.
\]
Thus $a-b=0$; the normalization condition then gives $a=b=1/2$.

Wave functions for states with $M_L<L$ are obtained from those found above
by applying the operator $\hat L_-$.
